\chapter{Khảo sát cửa hàng --- Metabase}

Chương này gồm \textbf{17 đường dẫn} còn lại, chia ba phần:

\begin{itemize}[leftmargin=*]
  \item \textbf{Phần A --- phía cửa hàng} (\rt{wujia_portal_inspection}, 8 đường dẫn): cửa hàng
        xem kết quả khảo sát và gửi khắc phục.
  \item \textbf{Phần B --- phía người đi khảo sát} (\rt{wujia_franchise_inspection}, 7 đường
        dẫn): màn nhân sự Ngô Gia dùng khi xuống cửa hàng chấm điểm. \textbf{Không phải màn
        portal của cửa hàng}, nhưng chạy trên cùng máy chủ và cùng đường dẫn web nên đưa vào
        đây cho đủ bản đồ.
  \item \textbf{Phần C --- Metabase} (\rt{wujia_metabase_connector}, 2 đường dẫn): nhúng bảng
        số liệu từ hệ thống báo cáo ngoài.
\end{itemize}

\textbf{Lưu ý về quyền sở hữu mã:} toàn bộ chương này là mã của anh Thái. Phiên tài liệu chỉ
\textbf{đọc để mô tả}, không sửa một dòng nào. Những điểm dev thấy cần bàn được ghi trong khung
vàng như mọi chương khác --- BA đọc để nắm, phần xử lý dev sẽ trao đổi trực tiếp với anh Thái.

\textbf{Ba điểm cần biết trước khi đọc:}

\begin{enumerate}[leftmargin=*]
  \item \textbf{Phần A phạm vi là MỌI cửa hàng tôi truy cập được} --- không phải cửa hàng đang
        chọn. Đây là cách làm giống Đổi trả, khác Đăng ký thi và Công nợ.
  \item \textbf{Điểm hiện trên portal được tính lại bằng công thức riêng của portal}, không lấy
        thẳng điểm đã chốt trong ERP --- xem mục Màn chi tiết. Hai nơi có thể ra hai con số.
  \item \textbf{Phần B và phần C chỉ kiểm ``đã đăng nhập'' ở phần lớn đường dẫn} --- chi tiết
        ghi trong từng mục.
\end{enumerate}

Số thật trên UAT: \uat{1} phiếu khảo sát (cửa hàng \emph{TP HCM Quận 1}, ngày 01/09, trạng thái
\emph{Chờ phản hồi}, điểm 86, xếp loại B), \uat{4} dòng tiêu chí trong đó \uat{2} không đạt,
\uat{0} tiêu chí đã được cửa hàng gửi khắc phục, \uat{0} dòng điểm danh, \uat{1} bảng Metabase.

% =============================================================================
\section{Màn Danh sách phiếu khảo sát}
\rt{/portal/inspection} · \rt{/portal/inspection/ajax}
\medskip

\mvao

Mục \emph{Khảo sát} trên thanh điều hướng. Đường dẫn \rt{/ajax} không phải màn riêng --- là
phần kết quả gọi ngầm khi đổi thẻ hoặc lọc. Tham số: \rt{tab}, \rt{date_from}, \rt{date_to},
\rt{search}, \rt{page}.

\mai

\begin{itemize}[leftmargin=*]
  \item Phải đăng nhập. \textbf{Không chặn vai trò.}
  \item Thấy phiếu của \textbf{mọi cửa hàng tôi truy cập được}.
  \item Không truy cập được cửa hàng nào thì \textbf{không thấy phiếu nào} --- điều kiện lọc cố
        ý để rỗng thay vì bỏ qua, nên không có kẽ hở xem nhầm cửa hàng khác.
  \item \textbf{Chỉ thấy phiếu đã chấm xong}: trạng thái \emph{Chờ phản hồi} hoặc
        \emph{Hoàn thành}. Phiếu người khảo sát đang làm dở (còn nháp) portal không thấy.
\end{itemize}

\mhien

\begin{hienbang}
Danh sách phiếu &
  \rt{wujia.franchise.inspection} &
  Phiếu thuộc cửa hàng của tôi, trạng thái \emph{Chờ phản hồi} hoặc \emph{Hoàn thành} &
  Ngày khảo sát mới nhất trước; \textbf{5} phiếu mỗi trang, \textbf{cố định} &
  \uat{1} phiếu · \uat{0} phiếu nháp \\ \addlinespace
Huy hiệu trạng thái &
  \rt{state} &
  Hai nhãn: \emph{Chờ phản hồi} (vàng) và \emph{Hoàn thành} (xanh) &
  --- &
  \uat{1} Chờ phản hồi \\ \addlinespace
Điểm và xếp loại &
  \rt{total_score} · \rt{wujia.franchise.inspection.grade} &
  Điểm tổng của phiếu; xếp loại lấy từ phiếu, nếu phiếu chưa gắn xếp loại thì
  \textbf{tra lại theo thang điểm} &
  --- &
  \uat{86} điểm, loại B \\ \addlinespace
Số tiêu chí đạt / không đạt &
  \rt{wujia.franchise.inspection.line} &
  Đếm các dòng tiêu chí; \emph{không đạt} = dòng bị đánh trượt &
  --- &
  \uat{2}/4 tiêu chí không đạt \\ \addlinespace
Nút \emph{Gửi khắc phục} &
  --- &
  Chỉ hiện khi phiếu đang \emph{Chờ phản hồi} và còn tiêu chí trượt \textbf{chưa gửi khắc
  phục}; nút trỏ thẳng tới tiêu chí đầu tiên còn thiếu &
  --- & --- \\ \addlinespace
Tên phiếu &
  \rt{name} &
  Phiếu đặt tên chung chung (\emph{``Khảo sát cửa hàng nhượng quyền''}) thì màn \textbf{tự đổi}
  thành \emph{``Lần dd/mm/yyyy''} cho dễ phân biệt &
  --- & --- \\
\end{hienbang}

\mloc

\begin{itemize}[leftmargin=*]
  \item \textbf{Ba thẻ}: \emph{Tất cả} · \emph{Chờ phản hồi} · \emph{Hoàn thành}.
  \item \textbf{Từ khoá} --- khớp một phần theo \textbf{tên phiếu}, \textbf{tên cửa hàng} hoặc
        \textbf{mã cửa hàng}.
  \item \textbf{Khoảng ngày} --- lọc theo \textbf{ngày khảo sát}; nhận cả dạng
        \emph{dd/mm/yyyy} lẫn \emph{yyyy-mm-dd}.
  \item \textbf{Không đổi được số dòng mỗi trang} --- cố định 5.
\end{itemize}

\mkiem{một phiếu}

Không kiểm gì ở danh sách.

\mghi

\textbf{Không ghi gì.}

\mhoi

\begin{ask}
\begin{enumerate}[leftmargin=*]
  \item \textbf{Gõ ngày ngược thì màn hiện rỗng, không báo gì} --- khác năm màn khác trong
        portal (Giao hàng, Báo cáo, Thông báo, Đăng ký thi, Lịch sử mua) vốn giữ hai ô đã gõ và
        báo tại thanh lọc. Anh có muốn đồng bộ không?
  \item \textbf{Cỡ trang cố định 5} --- nhỏ hơn hẳn các màn khác (10 hoặc 12) và không đổi được.
  \item \textbf{Phiếu còn nháp portal không thấy} --- đúng ý anh chứ? (Nghĩa là cửa hàng chỉ
        biết kết quả sau khi người khảo sát bấm hoàn tất.)
\end{enumerate}
\end{ask}

% =============================================================================
\section{Màn Chi tiết phiếu khảo sát}
\rt{/portal/inspection/detail/<int:inspection_id>}
\medskip

\mvao

Bấm một phiếu ở màn danh sách.

\mai

\begin{kiembang}
1 & Phiếu tồn tại & Về danh sách \\
2 & Trạng thái là \emph{Chờ phản hồi} hoặc \emph{Hoàn thành} & Về danh sách \\
3 & Phiếu thuộc \textbf{một trong các cửa hàng của tôi} & Về danh sách \\
\end{kiembang}

Cả ba trường hợp đều đưa về danh sách \textbf{không có thông báo}.

\mhien

\begin{hienbang}
Đầu phiếu &
  \rt{wujia.franchise.inspection} &
  Mã cửa hàng, ngày khảo sát, điểm tổng, xếp loại &
  --- &
  \uat{86}/100, loại B \\ \addlinespace
Bảng tóm tắt danh mục &
  \rt{wujia.franchise.inspection.line} &
  Mỗi danh mục một dòng. Danh mục thường hiện \emph{điểm đạt / điểm tối đa}; danh mục
  \textbf{vi phạm nghiêm trọng} hiện \textbf{số lỗi} thay cho điểm. &
  Theo thứ tự trong phiếu & --- \\ \addlinespace
Danh sách tiêu chí &
  \rt{wujia.franchise.inspection.line} &
  Nội dung tiêu chí, đạt hay không, điểm trừ, ghi chú của người khảo sát, ảnh hiện trường,
  ghi chú khắc phục và ảnh khắc phục của cửa hàng &
  Theo thứ tự trong phiếu &
  \uat{4} tiêu chí \\ \addlinespace
Nút \emph{Gửi khắc phục} &
  --- &
  Chỉ hiện khi phiếu đang \emph{Chờ phản hồi} và còn tiêu chí trượt chưa phản hồi &
  --- & --- \\
\end{hienbang}

\textbf{Điểm trên màn này được tính lại tại chỗ, không lấy điểm đã chốt trong ERP.} Cách tính
của portal: mỗi tiêu chí đạt được cộng số điểm bằng \textbf{điểm trừ của chính nó}; tiêu chí
\textbf{không khai điểm trừ thì được tính là 4 điểm}; và số điểm trừ lẻ bị \textbf{cắt phần
thập phân} (2,5 thành 2). Cột \emph{điểm tối đa} của mỗi danh mục cũng cộng theo cách đó, còn
\textbf{tổng điểm tối đa của phiếu luôn ghi cứng là 100}.

\textbf{Cách nhận biết ``vi phạm nghiêm trọng''} là \textbf{dò chữ}: hệ thống tìm cụm chữ
\emph{``nghiêm trọng''} trong tên danh mục hoặc trong nội dung tiêu chí. Đổi cách đặt tên danh
mục là cách nhận biết này hỏng.

\mloc

Không có ô lọc.

\mkiem{Gửi khắc phục}

Xem hai mục kế tiếp.

\mghi

Mở màn \textbf{không ghi gì}.

\mhoi

\begin{ask}
\begin{enumerate}[leftmargin=*]
  \item \textbf{Điểm trên portal và điểm trong ERP có thể lệch nhau} --- portal tự cộng lại
        theo quy ước riêng (tiêu chí không khai điểm trừ tính 4 điểm, điểm lẻ bị cắt phần thập
        phân), trong khi ERP dùng công thức chấm điểm của phiếu. Anh cần chốt: portal hiện đúng
        con số ERP, hay giữ cách tính riêng này?
  \item \textbf{Tổng điểm tối đa ghi cứng 100} --- bộ tiêu chí nào có thang khác 100 sẽ hiện
        sai tỉ lệ.
  \item \textbf{``Vi phạm nghiêm trọng'' nhận biết bằng cách dò chữ ``nghiêm trọng''} thay vì
        dùng cờ đánh dấu. Anh có muốn Ngô Gia đánh dấu hẳn một cờ trong danh mục không?
  \item \textbf{Có vài giá trị mẫu còn sót trong màn} --- phiếu thiếu mã cửa hàng, thiếu ngày
        hoặc thiếu điểm thì màn hiện giá trị mẫu (mã \emph{HN\_ST042}, ngày \emph{15/05/2024},
        điểm \emph{85}, loại \emph{B+}) thay vì để trống. Dev sẽ trao đổi với anh Thái, BA
        không cần lên task.
\end{enumerate}
\end{ask}

% =============================================================================
\section{Màn Gửi khắc phục cho một tiêu chí}
\rt{/portal/inspection/remediation/<int:line_id>} ·
\rt{/portal/remediation/<int:line_id>} · \rt{/portal/remediation}
\medskip

\mvao

Nút \emph{Gửi khắc phục} ở màn danh sách hoặc màn chi tiết. Ba đường dẫn cùng một màn:
đường dẫn đầy đủ, đường dẫn rút gọn, và \textbf{đường dẫn thiếu mã tiêu chí} --- gõ
\rt{/portal/remediation} không kèm số thì bị đưa về danh sách.

\mai

\begin{kiembang}
1 & Có mã tiêu chí & Về danh sách \\
2 & Tiêu chí tồn tại và thuộc một phiếu khảo sát có thật & Về danh sách \\
3 & Phiếu thuộc \textbf{một trong các cửa hàng của tôi} & Về danh sách \\
\end{kiembang}

\mhien

\begin{hienbang}
Thông tin tiêu chí &
  \rt{wujia.franchise.inspection.line} &
  Tên danh mục, mã tiêu chí, nội dung tiêu chí, \textbf{ghi chú của người khảo sát} &
  --- & --- \\ \addlinespace
Ảnh hiện trường &
  \rt{evidence_image} &
  Ảnh người khảo sát chụp lúc chấm --- tải qua đường dẫn ảnh chuẩn của hệ thống, có kiểm quyền
  theo cửa hàng &
  --- & --- \\ \addlinespace
Mốc thời gian &
  \rt{create_date} &
  Giờ và ngày \textbf{tạo phiếu khảo sát} --- không phải ngày khảo sát ghi trong phiếu &
  --- & --- \\
\end{hienbang}

\mloc

Hai ô: \textbf{ghi chú khắc phục} và \textbf{ảnh minh chứng} (chụp hoặc chọn từ máy).

\mkiem{Gửi}

Xem mục kế tiếp.

\mghi

Mở màn \textbf{không ghi gì}.

\mhoi

\begin{ask}
\textbf{Mốc thời gian trên màn là giờ tạo phiếu, không phải ngày khảo sát} --- hai con số này
lệch nhau khi người khảo sát tạo phiếu trước rồi vài hôm sau mới đi. Nên đổi sang ngày khảo sát.
\end{ask}

% =============================================================================
\section{Gửi nội dung khắc phục}
\rt{/portal/inspection/remediation/submit} · \rt{/portal/remediation/submit}
\medskip

\mvao

Không phải màn --- nút \emph{Gửi} ở màn khắc phục. Hai đường dẫn dẫn tới cùng một việc.

\mai

Chỉ kiểm \textbf{đã đăng nhập}. \textbf{Không} kiểm tiêu chí có thuộc cửa hàng của người gửi
hay không --- khác hẳn màn hiển thị ngay trên, vốn kiểm đủ ba lớp.

\mhien

Không hiện màn --- trả về câu báo thành công kèm đường dẫn quay lại màn chi tiết.

\mloc

Hai ô như mục trên. Ảnh nhận được ở \textbf{bốn dạng} khác nhau (tệp tải lên, chuỗi ảnh nhúng,
chuỗi mã hoá, dữ liệu nhị phân) --- viết vậy để cả bản điện thoại lẫn bản máy tính dùng chung
một đường dẫn.

\mkiem{Gửi}

\begin{kiembang}
1 & Có mã tiêu chí & ``Thiếu ID tiêu chí.'' \\
2 & Tiêu chí tồn tại & ``Tiêu chí không tồn tại.'' \\
3 & Tiêu chí \textbf{bắt buộc ảnh} thì phải có ảnh (ảnh mới hoặc ảnh đã gửi trước đó) &
    ``Tiêu chí này yêu cầu bắt buộc phải có ảnh minh chứng.'' \\
4 & Phải có \textbf{ít nhất} ghi chú hoặc ảnh &
    ``Vui lòng nhập ghi chú khắc phục hoặc tải lên ảnh minh chứng trước khi gửi.'' \\
\end{kiembang}

\mghi

\begin{itemize}[leftmargin=*]
  \item Ghi \textbf{ghi chú khắc phục} và \textbf{ảnh khắc phục} lên dòng tiêu chí.
  \item Đổi trạng thái khắc phục của \textbf{riêng dòng đó} sang \emph{Đã khắc phục}.
  \item \textbf{Không} đổi trạng thái của cả phiếu --- phiếu vẫn ở \emph{Chờ phản hồi} cho tới
        khi Ngô Gia tự chuyển. Cửa hàng gửi khắc phục đủ cả bốn tiêu chí thì phiếu vẫn hiện
        \emph{Chờ phản hồi}.
  \item \textbf{Không} báo cho ai --- không email, không việc cần làm.
\end{itemize}

\mhoi

\begin{ask}
\begin{enumerate}[leftmargin=*]
  \item \textbf{Đường dẫn gửi khắc phục không kiểm cửa hàng.} Người đăng nhập bất kỳ, nếu biết
        mã của một dòng tiêu chí, có thể ghi ghi chú và ảnh lên tiêu chí của cửa hàng khác.
        Đây là điểm dev sẽ trao đổi với anh Thái --- \textbf{BA không cần lên task}, nhưng cần
        biết vì nó ảnh hưởng tới câu ``ai xem được gì''.
  \item \textbf{Ảnh khắc phục không kiểm định dạng, không kiểm dung lượng} --- khác hẳn Đổi trả
        (chương 6) và Đăng ký thi (chương 8). Cùng nhóm với điểm trên.
  \item \textbf{Gửi khắc phục xong phiếu vẫn ``Chờ phản hồi''} và không ai được báo. Anh muốn
        phiếu tự chuyển trạng thái khi mọi tiêu chí trượt đã có phản hồi không?
  \item \textbf{Gửi rồi vẫn gửi lại được} --- lần gửi sau đè lên lần trước, không lưu lịch sử.
\end{enumerate}
\end{ask}

% =============================================================================
\section{Màn Làm khảo sát tại cửa hàng}
\rt{/franchise/inspection/do/<int:inspection_id>}
\medskip

\mvao

Nhân sự Ngô Gia mở từ phiếu khảo sát trong ERP khi xuống cửa hàng. \textbf{Không nằm trên thanh
điều hướng của portal cửa hàng.}

\mai

\begin{itemize}[leftmargin=*]
  \item Phải đăng nhập.
  \item Mở được nếu là \textbf{quản trị hệ thống}, là \textbf{người được giao khảo sát}, hoặc
        là \textbf{bất kỳ nhân sự nội bộ nào}.
  \item \textbf{Phiếu chưa giao cho ai thì không kiểm gì} --- bất kỳ tài khoản đăng nhập nào
        cũng mở được, kể cả tài khoản cửa hàng. Ghi lại ở khung vàng bên dưới.
\end{itemize}

\mhien

\begin{hienbang}
Danh sách tiêu chí &
  \rt{wujia.franchise.inspection.line} &
  Toàn bộ dòng của phiếu, chia theo danh mục, kèm điểm trừ và các cờ ``bắt buộc ghi chú'' /
  ``bắt buộc ảnh'' &
  Theo thứ tự trong phiếu &
  \uat{4} tiêu chí \\ \addlinespace
Kết quả lần khảo sát trước &
  \rt{previous_line_id} &
  Với mỗi tiêu chí: lần trước đạt hay trượt, ghi chú và ảnh lần trước, tên người khảo sát lần
  trước &
  --- & --- \\ \addlinespace
Bài kiểm tra nhân viên &
  \rt{exam_line_ids} &
  Câu hỏi kiểm tra kiến thức cho nhân viên được chọn &
  --- & --- \\ \addlinespace
Bảng điểm danh &
  \rt{...attendance.line} &
  Nhân sự có mặt lúc khảo sát &
  --- &
  \uat{0} dòng \\ \addlinespace
Báo cáo doanh thu 3 tháng &
  \rt{report_line_ids} &
  Doanh thu, doanh thu bình quân, tỉ lệ hàng hoá, doanh số qua ứng dụng &
  --- & --- \\ \addlinespace
Ngôn ngữ nội dung &
  --- &
  Nội dung tiêu chí lưu \textbf{song ngữ trong cùng một ô}; màn \textbf{tự tách} phần tiếng
  Trung hoặc tiếng Việt theo ngôn ngữ tài khoản &
  --- & --- \\
\end{hienbang}

\mloc

Chấm đạt / không đạt từng tiêu chí, nhập ghi chú, chụp ảnh hiện trường, trả lời bài kiểm tra,
điểm danh nhân sự, ký tên.

\mkiem{Lưu}

Xem mục kế tiếp.

\mghi

Mở màn \textbf{không ghi gì}.

\mhoi

\begin{ask}
\textbf{Phiếu chưa giao cho người khảo sát nào thì ai đăng nhập cũng mở được màn này}, kể cả
tài khoản cửa hàng --- và màn này hiện cả kết quả lần trước lẫn bài kiểm tra. Dev sẽ trao đổi
với anh Thái; BA chỉ cần biết để trả lời khi cửa hàng hỏi.
\end{ask}

% =============================================================================
\section{Lưu bài khảo sát}
\rt{/franchise/inspection/do/<int:inspection_id>/save}
\medskip

\mvao

Không phải màn --- nút \emph{Lưu tạm} và nút \emph{Hoàn tất} ở màn làm khảo sát.

\mai

Chỉ kiểm \textbf{đã đăng nhập} và phiếu chưa khoá. \textbf{Không} kiểm người gửi có phải người
được giao khảo sát hay không.

\mhien

Trả về điểm vừa tính lại, xếp loại, trạng thái phiếu và kết quả bài kiểm tra.

\mloc

Nhận cả gói: các dòng tiêu chí, bài kiểm tra, tên nhân viên được kiểm, thâm niên, tình trạng
mặt bằng, người xác nhận, bảng điểm danh, ảnh chữ ký, báo cáo doanh thu.

\mkiem{Hoàn tất}

\begin{kiembang}
1 & Phiếu tồn tại & ``Inspection sheet does not exist'' \\
2 & Phiếu \textbf{chưa hoàn tất, chưa bị huỷ} &
    ``Inspection sheet is already completed or cancelled and cannot be edited!'' \\
3 & Bấm \emph{Hoàn tất} thì phải có \textbf{tên nhân viên được kiểm} &
    ``Please enter staff name before saving!'' \\
\end{kiembang}

Bấm \emph{Lưu tạm} thì bỏ qua bước 3 và \textbf{không khoá bài kiểm tra}; bấm \emph{Hoàn tất}
thì bài kiểm tra bị chốt, không sửa lại được.

\mghi

\begin{itemize}[leftmargin=*]
  \item Ghi từng dòng tiêu chí --- \textbf{chỉ ghi dòng thật sự thay đổi}, so sánh trước khi
        ghi để không đụng dữ liệu thừa.
  \item Chấm bài kiểm tra và ghi điểm từng câu.
  \item Ghi bảng điểm danh; nhân sự mới trong bảng điểm danh \textbf{được tạo thành thành viên
        cửa hàng} luôn.
  \item Ghi lại bảng doanh thu 3 tháng --- dòng không còn gửi lên thì \textbf{bị xoá}.
  \item Tính lại điểm phần tiêu chí, điểm bài kiểm tra, điểm tổng và xếp loại.
  \item \textbf{Có tiêu chí trượt thì phiếu tự chuyển từ nháp sang \emph{Chờ phản hồi}} ---
        đây chính là lúc phiếu hiện ra trên portal cửa hàng.
\end{itemize}

\mhoi

\begin{ask}
\begin{enumerate}[leftmargin=*]
  \item \textbf{Đường dẫn lưu không kiểm người gửi.} Bất kỳ tài khoản đăng nhập nào biết mã
        phiếu đều ghi được điểm, ghi chú, chữ ký và bảng điểm danh lên phiếu đang mở. Cùng nhóm
        với hai điểm trên --- dev trao đổi với anh Thái.
  \item \textbf{Điểm danh tự tạo thành viên cửa hàng.} Gõ nhầm tên là sinh một thành viên mới
        trong hồ sơ cửa hàng. Anh có muốn bước xác nhận trước khi tạo không?
  \item \textbf{Phiếu tự chuyển sang Chờ phản hồi ngay khi lưu có tiêu chí trượt} --- kể cả khi
        người khảo sát mới chỉ \emph{lưu tạm}. Cửa hàng sẽ thấy phiếu trước khi người khảo sát
        làm xong.
\end{enumerate}
\end{ask}

% =============================================================================
\section{Kéo doanh thu từ PosApp}
\rt{/franchise/inspection/do/<int:inspection_id>/sync_posapp}
\medskip

\mvao

Không phải màn --- nút đồng bộ trong phần báo cáo doanh thu ở màn làm khảo sát.

\mai

Chỉ kiểm đã đăng nhập và phiếu chưa khoá.

\mhien

Trả về ba dòng doanh thu vừa kéo về để màn vẽ lại bảng.

\mloc

Không có ô nhập --- bấm là chạy.

\mkiem{Đồng bộ}

\begin{kiembang}
1 & Phiếu tồn tại & ``Inspection sheet does not exist'' \\
2 & Phiếu chưa hoàn tất, chưa bị huỷ &
    ``Inspection sheet is already completed or cancelled!'' \\
3 & Gọi được sang PosApp & Trả về nguyên văn câu lỗi kỹ thuật của hệ thống \\
\end{kiembang}

\mghi

\textbf{Xoá sạch bảng doanh thu cũ rồi ghi lại bảng mới} --- số nhân sự gõ tay trước đó mất
hết, không hỏi lại.

\mhoi

\begin{ask}
\begin{enumerate}[leftmargin=*]
  \item \textbf{Bấm đồng bộ là xoá sạch số đã gõ tay}, không có bước xác nhận. Anh có muốn hỏi
        lại trước khi ghi đè không?
  \item \textbf{Lỗi kết nối hiện nguyên văn câu lỗi kỹ thuật} --- cùng nhóm chuẩn hoá thông báo
        lỗi mà dev đang làm.
\end{enumerate}
\end{ask}

% =============================================================================
\section{Bảng điểm danh nhân sự}
\rt{.../attendance/add} · \rt{.../attendance/save_member} ·
\rt{.../attendance/deactivate_member} · \rt{.../attendance/delete_line}
\medskip

\mvao

Không phải màn --- bốn nút trong bảng điểm danh ở màn làm khảo sát: thêm dòng, lưu dòng thành
thành viên cửa hàng, ngưng hiệu lực thành viên, xoá dòng.

\mai

Chỉ kiểm \textbf{đã đăng nhập}. Ba đường dẫn sau còn kiểm dòng điểm danh \textbf{đúng thuộc
phiếu ghi trên đường dẫn} --- nên không mượn được để sửa dòng của phiếu khác.

\mhien

Mỗi lần gọi đều trả về \textbf{bảng nhân sự của cửa hàng} và \textbf{số người có mặt} để màn
cập nhật ngay.

\mloc

Họ tên (bắt buộc khi thêm), vai trò, số điện thoại, ghi chú, có mặt hay không.

\mkiem{Thêm nhân sự}

\begin{kiembang}
1 & Phiếu tồn tại và chưa hoàn tất / chưa bị huỷ &
    ``Inspection survey is locked or does not exist!'' \\
2 & Có họ tên & ``Please enter staff full name!'' \\
3 & (ba nút còn lại) Dòng điểm danh tồn tại và thuộc đúng phiếu &
    ``Attendance line does not exist!'' \\
\end{kiembang}

\mghi

\begin{itemize}[leftmargin=*]
  \item \textbf{Thêm}: tạo dòng điểm danh, rồi \textbf{tự tạo thành viên cửa hàng} cho người đó
        --- tạo không được thì bỏ qua, dòng điểm danh vẫn còn.
  \item \textbf{Lưu thành thành viên}: sửa tên, vai trò, điện thoại trên dòng rồi tạo hoặc cập
        nhật thành viên cửa hàng tương ứng.
  \item \textbf{Ngưng hiệu lực}: tắt hiệu lực thành viên cửa hàng --- \textbf{đây là thao tác
        sửa hồ sơ cửa hàng từ màn khảo sát}.
  \item \textbf{Xoá dòng}: xoá hẳn dòng điểm danh; thành viên cửa hàng đã tạo \textbf{vẫn còn}.
\end{itemize}

\mhoi

\begin{ask}
\begin{enumerate}[leftmargin=*]
  \item \textbf{Màn khảo sát sửa được hồ sơ nhân sự của cửa hàng} --- thêm thành viên, ngưng
        hiệu lực thành viên. Anh xác nhận cho phép; nếu không thì cần tách quyền.
  \item \textbf{Xoá dòng điểm danh không gỡ thành viên đã tạo} --- gõ nhầm rồi xoá dòng thì
        người đó vẫn nằm trong danh sách nhân sự cửa hàng.
\end{enumerate}
\end{ask}

% =============================================================================
\section{Bảng số liệu Metabase}
\rt{/metabase/embed/<int:dashboard_id>} · \rt{/metabase/embed_url/<int:dashboard_id>}
\medskip

\mvao

Không nằm trên thanh điều hướng portal. Đường dẫn thứ nhất mở trang có khung nhúng; đường dẫn
thứ hai chỉ trả về địa chỉ nhúng cho màn khác tự gắn vào.

\mai

\begin{kiembang}
1 & Đã đăng nhập & Trang đăng nhập \\
2 & Bảng tồn tại và đang bật & Báo không tìm thấy \\
3 & Bảng cho phép \textbf{công ty} hiện tại &
    Trang báo ``Access Denied'' kèm tên công ty \\
4 & \textbf{Nếu bảng có khai nhóm quyền}: tôi thuộc một trong các nhóm đó (hoặc là quản trị) &
    Trang báo ``Access Denied'' \\
5 & Có khai máy chủ Metabase và máy chủ đang bật & Trang báo ``Configuration Error'' \\
6 & Máy chủ có thư viện ký chữ ký số & Trang báo ``System Error'' \\
\end{kiembang}

\textbf{Bước 4 chỉ chạy khi bảng có khai nhóm quyền.} Bảng \textbf{không khai nhóm nào} thì bỏ
qua hoàn toàn --- mọi tài khoản đăng nhập đều xem được. Trên UAT, bảng duy nhất
(\emph{``Báo cáo nhượng quyền''}) \textbf{không khai nhóm nào}.

\mhien

\begin{hienbang}
Khung nhúng &
  \rt{metabase.dashboard} &
  Bảng số liệu nằm trên máy chủ Metabase ngoài, gọi qua một chữ ký số có hạn &
  --- &
  \uat{1} bảng, đang bật \\ \addlinespace
Hạn của chữ ký &
  \rt{token_expiry_minutes} &
  Mặc định \textbf{10 phút}; hết hạn thì khung nhúng báo lỗi, phải tải lại trang &
  --- & --- \\ \addlinespace
Chiều cao khung &
  \rt{height} &
  Theo cấu hình của bảng, mặc định 800 &
  --- & --- \\
\end{hienbang}

\mloc

Không có ô lọc --- mọi thao tác lọc diễn ra bên trong khung Metabase, hệ thống này không can
thiệp. \textbf{Không truyền tham số lọc nào sang Metabase} --- nghĩa là bảng số liệu
\textbf{không tự giới hạn theo cửa hàng của người xem}.

\mkiem{mở bảng}

Xem bảng ở mục \emph{Ai xem được gì}.

\mghi

\textbf{Không ghi gì.}

\mhoi

\begin{ask}
\begin{enumerate}[leftmargin=*]
  \item \textbf{Bảng không khai nhóm quyền thì ai đăng nhập cũng xem được} --- kể cả tài khoản
        cửa hàng. Bảng duy nhất trên UAT đang ở trạng thái đó. Nếu định mở Metabase cho cửa
        hàng thì phải khai nhóm trước.
  \item \textbf{Không truyền cửa hàng sang Metabase} --- người xem thấy số liệu toàn hệ, không
        phải số liệu cửa hàng mình. Anh xác nhận đây là bảng dành cho nội bộ Ngô Gia.
\end{enumerate}
\end{ask}
